
\documentclass[12pt,a4paper]{article}

% ქართული შრიფტი
\usepackage{fontspec}
\setmainfont{DejaVu Sans}

% ბიბლიოთეკები
\usepackage{geometry}
\usepackage{setspace}
\usepackage{xcolor}
\usepackage{titlesec}
\usepackage{enumitem}
\usepackage{tikz}
\usepackage{fancyhdr}
\usepackage{tcolorbox}
\usepackage{array}
\usepackage{graphicx}

\usetikzlibrary{
    arrows.meta,
    positioning,
    shapes.geometric,
    calc,
    decorations.pathmorphing
}

% გვერდის დიზაინი
\geometry{
    left=2.2cm,
    right=2.2cm,
    top=2.3cm,
    bottom=2.3cm
}

\onehalfspacing
\setlength{\parindent}{0.7cm}

% ფერები
\definecolor{navy}{RGB}{22,55,92}
\definecolor{blue}{RGB}{40,120,190}
\definecolor{lightblue}{RGB}{225,241,252}
\definecolor{cyan}{RGB}{30,170,190}
\definecolor{orange}{RGB}{240,145,50}
\definecolor{lightorange}{RGB}{255,240,220}
\definecolor{green}{RGB}{55,145,105}
\definecolor{lightgreen}{RGB}{225,245,235}
\definecolor{red}{RGB}{190,65,65}
\definecolor{lightgray}{RGB}{242,245,248}

% სათაურები
\titleformat{\section}
{\Large\bfseries\color{navy}}
{\thesection.}{0.5em}{}

\titleformat{\subsection}
{\large\bfseries\color{blue}}
{\thesubsection.}{0.5em}{}

% გვერდის ნომრები
\pagestyle{fancy}
\fancyhf{}
\fancyhead[L]{\small\color{navy} ხელოვნური კუნთი}
\fancyhead[R]{\small Artificial Muscle}
\fancyfoot[C]{\thepage}
\renewcommand{\headrulewidth}{0.4pt}

% ინფორმაციის ბლოკები
\newtcolorbox{infobox}{
    colback=lightblue,
    colframe=blue,
    arc=3mm,
    boxrule=0.8pt,
    left=3mm,
    right=3mm,
    top=2mm,
    bottom=2mm
}

\newtcolorbox{examplebox}{
    colback=lightorange,
    colframe=orange,
    arc=3mm,
    boxrule=0.8pt,
    left=3mm,
    right=3mm,
    top=2mm,
    bottom=2mm
}

\begin{document}

% =========================
% ყდის გვერდი
% =========================

\begin{titlepage}
\begin{center}

\vspace*{1.5cm}

{\Huge\bfseries\color{navy} ხელოვნური კუნთი}

\vspace{0.5cm}

{\LARGE\color{blue} ARTIFICIAL MUSCLE}

\vspace{0.8cm}

{\large ტექნოლოგია, რომელიც მექანიკურ სამყაროს მოძრაობას სძენს}

\vspace{1.5cm}

% დეკორატიული ილუსტრაცია
\begin{tikzpicture}[scale=1.1]

% რობოტული მკლავი
\draw[line width=15pt,blue!60,line cap=round]
(0,0)--(1.7,1.3);

\draw[line width=13pt,cyan!65,line cap=round]
(1.7,1.3)--(3.2,2.2);

\fill[orange] (1.7,1.3) circle (0.18);
\fill[navy] (0,0) circle (0.22);
\fill[navy] (3.2,2.2) circle (0.15);

% თითები
\draw[line width=4pt,blue!60,line cap=round]
(3.2,2.2)--(3.8,2.7);
\draw[line width=4pt,blue!60,line cap=round]
(3.2,2.2)--(4.0,2.3);
\draw[line width=4pt,blue!60,line cap=round]
(3.2,2.2)--(3.8,1.9);

% კუნთის ხაზები
\draw[red,line width=2pt]
(0.5,0.4)--(1.5,1.1);
\draw[red,line width=2pt]
(1.9,1.5)--(2.9,2.0);

\node[below] at (1.6,-0.5)
{\small რობოტული მოძრაობის მოდელი};

\end{tikzpicture}

\vspace{1.5cm}

\begin{infobox}
\centering
\textbf{პროექტის მთავარი მიზანი}

ხელოვნური კუნთების მუშაობის პრინციპების,
სახეობების და თანამედროვე ტექნოლოგიებში
მათი გამოყენების შესწავლა.
\end{infobox}

\vfill

{\large სასწავლო პროექტი}

\vspace{0.5cm}

\end{center}
\end{titlepage}

% =========================
% შესავალი
% =========================

\section{რა არის ხელოვნური კუნთი?}

ხელოვნური კუნთი არის სპეციალური მოწყობილობა ან მასალა, რომელსაც შეუძლია ადამიანის ბუნებრივი კუნთის მსგავსი მოძრაობის შესრულება. მას შეუძლია შეკუმშვა, გაწელვა, მოხრა და მექანიკური ძალის წარმოქმნა.

ბუნებრივი კუნთი მოძრაობს ნერვული სისტემიდან მიღებული სიგნალების დახმარებით. ხელოვნური კუნთი კი ენერგიის სხვადასხვა წყაროს იყენებს, მაგალითად, ელექტროენერგიას, შეკუმშულ ჰაერს, სითხის წნევას ან ტემპერატურის ცვლილებას.

\begin{infobox}
\textbf{მთავარი იდეა:} ხელოვნური კუნთი ენერგიას გარდაქმნის მექანიკურ მოძრაობად. მისი დახმარებით შესაძლებელია რობოტის თითის მოხრა, მექანიკური ხელის ამოძრავება ან პროთეზის მოძრაობა.
\end{infobox}

\section{როგორ მუშაობს ხელოვნური კუნთი?}

მისი მუშაობის პრინციპი დამოკიდებულია გამოყენებულ ტექნოლოგიაზე. თუმცა, ზოგადად, პროცესი რამდენიმე ეტაპს მოიცავს:

\begin{center}
\begin{tikzpicture}[
node distance=0.65cm,
box/.style={
draw=blue,
fill=lightblue,
rounded corners=3mm,
minimum width=3cm,
minimum height=1cm,
align=center,
font=\small\bfseries
},
arr/.style={-{Stealth},very thick,navy}
]

\node[box] (a) {ენერგიის მიღება};
\node[box,right=of a] (b) {ენერგიის გარდაქმნა};
\node[box,right=of b] (c) {კუნთის მოძრაობა};
\node[box,right=of c] (d) {მექანიზმის ამოძრავება};

\draw[arr] (a)--(b);
\draw[arr] (b)--(c);
\draw[arr] (c)--(d);

\end{tikzpicture}
\end{center}

მაგალითად, როდესაც ხელოვნური კუნთი იღებს ელექტრულ სიგნალს, ის შეიძლება შეიკუმშოს ან შეიცვალოს მისი ფორმა. ეს ცვლილება წარმოქმნის ძალას, რომელსაც შეუძლია სხვა მექანიზმის ამოძრავება.

% =========================
% პნევმატიკური კუნთი
% =========================

\section{პნევმატიკური ხელოვნური კუნთი}

პნევმატიკური კუნთი მუშაობს შეკუმშული ჰაერის დახმარებით. მისი ერთ-ერთი ცნობილი სახეობაა მაკკიბენის კუნთი (McKibben Muscle).

ის შედგება მოქნილი შიდა მილისა და გარე ბადისებრი გარსისგან.

\subsection{როგორ ხდება შეკუმშვა?}

როდესაც შიდა მილში შეკუმშული ჰაერი შედის, მილი სიგანეში ფართოვდება. გარე ბადისებრი გარსი კი მის სიგრძეში გაფართოებას ზღუდავს. შედეგად, მთელი კონსტრუქცია სიგრძეში მოკლდება და წარმოქმნის გამწევ ძალას.

\begin{center}
\textbf{\color{navy} პნევმატიკური კუნთის მუშაობის სქემა}

\vspace{0.5cm}

\begin{tikzpicture}[scale=0.95]

% პირველი მდგომარეობა
\node[font=\bfseries,color=blue] at (2,3.2)
{ჰაერის გარეშე};

\draw[fill=lightblue,draw=blue,thick,rounded corners=5pt]
(0,1.3) rectangle (4,2.1);

\foreach \x in {0.2,0.6,...,3.8}{
\draw[blue!70,thick] (\x,1.3)--(\x+0.35,2.1);
}
\foreach \x in {0.2,0.6,...,3.8}{
\draw[blue!70,thick] (\x,2.1)--(\x+0.35,1.3);
}

\node[below] at (2,0.9) {\small გრძელი მდგომარეობა};

% ისარი
\draw[-{Stealth},very thick,orange]
(4.5,1.7)--(5.8,1.7);

\node[above] at (5.15,1.8)
{\scriptsize ჰაერი};

% მეორე მდგომარეობა
\node[font=\bfseries,color=red] at (8,3.2)
{შეკუმშული};

\draw[fill=lightorange,draw=red,thick,rounded corners=5pt]
(6.5,1.3) rectangle (9.5,2.1);

\foreach \x in {6.6,7.1,...,9.2}{
\draw[red!80,thick] (\x,1.3)--(\x+0.5,2.1);
}
\foreach \x in {6.6,7.1,...,9.2}{
\draw[red!80,thick] (\x,2.1)--(\x+0.5,1.3);
}

\node[below] at (8,0.9) {\small მოკლე მდგომარეობა};

\draw[-{Stealth},red,thick]
(6.5,2.5)--(7.1,2.5);
\draw[-{Stealth},red,thick]
(9.5,2.5)--(8.9,2.5);

\end{tikzpicture}
\end{center}

\begin{examplebox}
\textbf{რეალური მაგალითი:}

ასეთი კუნთი შეიძლება გამოიყენონ რობოტულ ხელში. როდესაც ის იკუმშება, თითი იხრება, ხოლო როდესაც ჰაერი გამოიდევნება, თითი საწყის მდგომარეობას უბრუნდება.
\end{examplebox}

% =========================
% რობოტული ხელი
% =========================

\section{ხელოვნური კუნთები რობოტულ ხელში}

რობოტული ხელი ერთ-ერთი ყველაზე საინტერესო მაგალითია, სადაც ხელოვნური კუნთების გამოყენება შეიძლება.

ადამიანის ხელში კუნთები და მყესები თითების მოძრაობას უზრუნველყოფს. რობოტულ ხელში კი ამ ფუნქციის შესრულება შეუძლიათ ხელოვნურ აქტუატორებს, რომლებიც თითებს ხრიან და შლიან.

\begin{center}
\begin{tikzpicture}[scale=1.0]

% ხელისგული
\draw[fill=lightblue,draw=blue,line width=1.5pt]
(1.2,0.2) rectangle (3.1,2.3);

% თითები
\foreach \x/\h in {1.35/1.8,1.8/2.2,2.3/2.0,2.75/1.6}{
\draw[fill=lightblue,draw=blue,line width=1.5pt,rounded corners=3pt]
(\x,2.3) rectangle (\x+0.3,2.3+\h);
\draw[orange,line width=3pt]
(\x+0.15,2.5)--(\x+0.15,2.3+\h-0.3);
}

% ცერი
\draw[fill=lightblue,draw=blue,line width=1.5pt,rounded corners=3pt]
(0.3,0.6) rectangle (1.2,1.0);
\draw[orange,line width=3pt]
(0.55,0.8)--(1.0,0.8);

% კუნთის ხაზები
\draw[red,line width=2pt]
(1.5,0.4)--(1.5,1.9);
\draw[red,line width=2pt]
(2.0,0.4)--(2.0,2.0);
\draw[red,line width=2pt]
(2.5,0.4)--(2.5,1.8);

\node[red,right] at (3.3,1.5)
{\small ხელოვნური კუნთები};

\node[blue,right] at (3.3,2.3)
{\small მოძრავი თითები};

\node[orange,right] at (3.3,0.7)
{\small გამწევი ელემენტები};

\end{tikzpicture}
\end{center}

\textit{სურათი არის გამარტივებული სასწავლო სქემა და არა კონკრეტული რობოტის ზუსტი კონსტრუქცია.}

\subsection{როგორ მოძრაობს რობოტის თითი?}

\begin{enumerate}
    \item მართვის სისტემა აგზავნის სიგნალს.
    \item აქტუატორი წარმოქმნის ძალას.
    \item ძალა გადაეცემა თითის მექანიზმს.
    \item თითი იხრება ან იშლება.
\end{enumerate}

ეს ტექნოლოგია განსაკუთრებით საინტერესოა პროთეზებისა და დამხმარე რობოტების შესაქმნელად.

% =========================
% ტიპები
% =========================

\section{ხელოვნური კუნთების სხვადასხვა ტიპები}

ხელოვნური კუნთების რამდენიმე ძირითადი ტექნოლოგია არსებობს.

\subsection{1. პნევმატიკური კუნთი}

\begin{center}
\begin{tikzpicture}
\draw[fill=lightblue,draw=blue,thick,rounded corners=5pt]
(0,0) rectangle (4,0.8);
\foreach \x in {0.2,0.7,...,3.7}{
\draw[blue,thick] (\x,0)--(\x+0.5,0.8);
\draw[blue,thick] (\x,0.8)--(\x+0.5,0);
}
\node[below] at (2,-0.2) {\small შეკუმშული ჰაერი};
\end{tikzpicture}
\end{center}

იყენებს ჰაერის წნევას და შეუძლია შეკუმშვა. გამოიყენება რობოტულ ხელებსა და მოქნილ მექანიზმებში.

\subsection{2. ელექტროაქტიური პოლიმერი}

\begin{center}
\begin{tikzpicture}
\draw[fill=lightorange,draw=orange,thick]
(0,0) rectangle (3,0.3);
\draw[fill=lightblue,draw=blue,thick]
(0,0.4) rectangle (3,0.7);
\draw[-{Stealth},red,thick] (1.5,1.3)--(1.5,0.8);
\node[above] at (1.5,1.35) {\small ელექტრული სიგნალი};
\node[below] at (1.5,-0.2) {\small მოქნილი პოლიმერული ფენა};
\end{tikzpicture}
\end{center}

ელექტრული ველის ზემოქმედებით ზოგიერთი პოლიმერი ფორმას ან ზომას იცვლის. ეს მასალები განსაკუთრებით საინტერესოა მცირე ზომის მოქნილი მოწყობილობებისთვის.

\subsection{3. ფორმის მეხსიერების შენადნობი}

\begin{center}
\begin{tikzpicture}
\draw[gray,very thick,decorate,
decoration={coil,segment length=4pt,amplitude=5pt}]
(0,0)--(2.5,0);
\draw[-{Stealth},orange,thick] (1.2,0.9)--(1.2,0.3);
\node[above] at (1.2,0.9) {\small გათბობა};
\node[below] at (1.2,-0.3) {\small ფორმის შეცვლა};
\end{tikzpicture}
\end{center}

ეს შენადნობები ტემპერატურის ცვლილებისას წინასწარ განსაზღვრულ ფორმას უბრუნდებიან. ამის გამოყენება შესაძლებელია პატარა მექანიზმებში.

\subsection{4. ჰიდრავლიკური აქტუატორი}

\begin{center}
\begin{tikzpicture}
\draw[fill=lightblue,draw=blue,thick]
(0,0) rectangle (2.8,0.8);
\draw[fill=orange,draw=orange]
(2.8,0.25) rectangle (3.8,0.55);
\draw[-{Stealth},blue,thick] (-0.7,0.4)--(0,0.4);
\node[above] at (1.4,0.9) {\small სითხის წნევა};
\node[below] at (1.4,-0.2) {\small ჰიდრავლიკური ცილინდრი};
\end{tikzpicture}
\end{center}

სითხის წნევის დახმარებით ქმნის მექანიკურ ძალას. გამოიყენება ისეთ სისტემებში, რომლებსაც ძლიერი მოძრაობა სჭირდებათ.

% =========================
% გამოყენება
% =========================

\section{სად გამოიყენება ხელოვნური კუნთები?}

\subsection{მედიცინა და პროთეზები}

\begin{center}
\begin{tikzpicture}
\draw[fill=lightblue,draw=blue,thick,rounded corners=4pt]
(0,0) rectangle (1.2,2);
\draw[orange,line width=4pt] (0.6,0.4)--(0.6,1.6);
\draw[blue,line width=2pt] (0.3,0.5)--(0.9,1.5);
\node[below] at (0.6,-0.2) {\small პროთეზი};
\end{tikzpicture}
\end{center}

ხელოვნური აქტუატორები შეიძლება გამოყენებული იყოს პროთეზებში, რათა მათ სხვადასხვა მოძრაობის შესრულება შეძლონ.

\subsection{რობოტიკა}

\begin{center}
\begin{tikzpicture}
% რობოტის სხეული
\draw[fill=lightblue,draw=blue,thick,rounded corners=5pt]
(0,0) rectangle (1.5,1.5);
\draw[fill=navy] (0.3,1.1) circle (0.1);
\draw[fill=navy] (1.2,1.1) circle (0.1);
\draw[blue,line width=5pt] (0.3,0)--(-0.3,-0.8);
\draw[blue,line width=5pt] (1.2,0)--(2,-0.8);
\draw[orange,line width=3pt] (1.5,0.7)--(2.1,0.3);
\node[below] at (0.7,-1) {\small რობოტი};
\end{tikzpicture}
\end{center}

რობოტებში ხელოვნური კუნთები ხელს უწყობს მოქნილი მოძრაობების შესრულებას და საგნებთან ურთიერთქმედებას.

\subsection{ეგზოჩონჩხები}

ეგზოჩონჩხი არის სპეციალური მოწყობილობა, რომელიც ადამიანის სხეულზე მაგრდება და მოძრაობის შესრულებაში ეხმარება. მასში შეიძლება გამოყენებული იყოს ელექტრული, პნევმატიკური ან ჰიდრავლიკური აქტუატორები.

\subsection{სარეაბილიტაციო მოწყობილობები}

ასეთი ტექნოლოგიები შეიძლება გამოიყენებოდეს მოძრაობის დამხმარე მოწყობილობებში, რომლებიც ადამიანს გარკვეული მოძრაობების შესრულებაში ეხმარება.

% =========================
% შედარება
% =========================

\section{უპირატესობები და ნაკლოვანებები}

\begin{center}
\renewcommand{\arraystretch}{1.7}
\begin{tabular}{|p{6cm}|p{6cm}|}
\hline
\rowcolor{lightgreen}
\textbf{უპირატესობები} &
\textbf{ნაკლოვანებები} \\
\hline
მოქნილი მოძრაობის შესაძლებლობა &
ზოგიერთ სისტემას რთული მართვა სჭირდება \\
\hline
ზოგიერთი ტიპის მცირე წონა &
ენერგიის მუდმივი მიწოდების საჭიროება \\
\hline
რობოტებთან ინტეგრირების შესაძლებლობა &
მასალების ცვეთა და შეზღუდული გამძლეობა \\
\hline
მედიცინაში გამოყენების პოტენციალი &
წარმოების მაღალი ღირებულება ზოგიერთ შემთხვევაში \\
\hline
\end{tabular}
\end{center}

% =========================
% მომავალი
% =========================

\section{ხელოვნური კუნთების მომავალი}

ხელოვნური კუნთების ტექნოლოგია სწრაფად ვითარდება. მეცნიერები ცდილობენ შექმნან მასალები, რომლებიც უფრო გამძლე, მსუბუქი და ეფექტური იქნება.

მომავალში მათი გამოყენება შესაძლებელია:

\begin{itemize}
    \item უფრო მოქნილ ჰუმანოიდ რობოტებში;
    \item განვითარებულ სამედიცინო პროთეზებში;
    \item ადამიანის მოძრაობის დამხმარე სისტემებში;
    \item მცირე ზომის სამედიცინო რობოტებში;
    \item თანამედროვე სამრეწველო მექანიზმებში.
\end{itemize}

\begin{infobox}
\textbf{საინტერესო ფაქტი:}

ხელოვნური კუნთების მიზანი მხოლოდ რობოტის მოძრაობა არ არის. მათი დახმარებით მეცნიერები ცდილობენ შექმნან ისეთი სისტემები, რომლებიც ბუნებრივი კუნთების მოქნილობასა და მოძრაობის პრინციპებს მიუახლოვდება.
\end{infobox}

% =========================
% დასკვნა
% =========================

\section{დასკვნა}

ხელოვნური კუნთი თანამედროვე მეცნიერებისა და ინჟინერიის ერთ-ერთი მნიშვნელოვანი მიმართულებაა. ის აერთიანებს ფიზიკის, ქიმიის, მასალათმცოდნეობისა და რობოტიკის პრინციპებს.

მისი მუშაობის შესწავლა გვიჩვენებს, თუ როგორ შეიძლება სხვადასხვა ენერგიის მექანიკურ მოძრაობად გარდაქმნა. ეს ტექნოლოგია უკვე გამოიყენება სხვადასხვა საინჟინრო და სამედიცინო მიმართულებაში და მომავალში კიდევ უფრო მნიშვნელოვანი გახდება.

ხელოვნური კუნთების განვითარება მეცნიერებს საშუალებას მისცემს შექმნან უფრო მოქნილი, ზუსტი და ადამიანის მოძრაობებთან მიახლოებული რობოტები.

\vspace{0.5cm}

\begin{center}
\begin{tcolorbox}[
colback=navy,
colframe=navy,
width=0.9\textwidth,
arc=3mm
]
\centering
\color{white}
\Large\bfseries
ხელოვნური კუნთი

\vspace{0.2cm}

\normalsize
მომავლის ტექნოლოგია, რომელიც ენერგიას მოძრაობად გარდაქმნის.
\end{tcolorbox}
\end{center}

\end{document}
